\chapter{Pagination}

\section{The skip/limit Formula}

Returning every matching document in one response doesn't scale once a user
has hundreds of tasks. Pagination splits results into pages using two
numbers derived from the requested page: how many documents to skip, and how
many to return.

\begin{lstlisting}[language=JavaScript]
const page = Number(req.query.page) || 1;
const limit = Number(req.query.limit) || 10;
const skip = (page - 1) * limit;
\end{lstlisting}

For page 1 with a limit of 10, \texttt{skip} is 0 --- return the first 10
documents. For page 2, \texttt{skip} is 10 --- skip the first 10 and return
the next 10. And so on.

\begin{diagrambox}[Pagination Flow]
\begin{verbatim}
Frontend requests page=2, limit=10
        |
        v
   GET /api/tasks?page=2&limit=10
        |
        v
   skip = (2 - 1) * 10 = 10
        |
        v
   MongoDB: Task.find(filter).skip(10).limit(10)
        |
        v
   Response: { data, page: 2, limit: 10, total, totalPages }
\end{verbatim}
\end{diagrambox}

\section{Backend Implementation}

Getting the total document count alongside the page of results lets the
frontend render page numbers and disable ``Next'' on the last page.

\begin{lstlisting}[language=JavaScript]
// backend/src/controllers/task.controller.js
const Task = require("../models/Task");

const getTasks = async (req, res, next) => {
  try {
    const page = Number(req.query.page) || 1;
    const limit = Number(req.query.limit) || 10;
    const skip = (page - 1) * limit;

    const filter = { owner: req.user.id };
    if (req.query.status) filter.status = req.query.status;
    if (req.query.search) {
      filter.title = { $regex: req.query.search, $options: "i" };
    }

    const sortObj = { createdAt: -1 };

    const [tasks, total] = await Promise.all([
      Task.find(filter).sort(sortObj).skip(skip).limit(limit),
      Task.countDocuments(filter),
    ]);

    const totalPages = Math.ceil(total / limit);

    res.status(200).json({
      data: tasks,
      page,
      limit,
      total,
      totalPages,
    });
  } catch (error) {
    next(error);
  }
};

module.exports = { getTasks };
\end{lstlisting}

\begin{notebox}[NOTE]
\texttt{Task.countDocuments(filter)} must use the \emph{same} filter object
as the \texttt{find()} call, otherwise \texttt{total} and \texttt{totalPages}
won't correspond to the actual filtered result set.
\end{notebox}

\subsection{Response Shape}

This is the exact shape referenced by the API reference table in
Chapter~8 --- every paginated list endpoint in the Task Management System
follows it consistently.

\begin{lstlisting}[language=JavaScript]
{
  "data": [ /* array of up to `limit` task objects */ ],
  "page": 2,
  "limit": 10,
  "total": 45,
  "totalPages": 5
}
\end{lstlisting}

\section{Frontend: Pagination Controls}

\begin{lstlisting}[language=JavaScript]
// frontend/src/components/Pagination.jsx
function Pagination({ page, totalPages, onPageChange }) {
  const pageNumbers = Array.from({ length: totalPages }, (_, i) => i + 1);

  return (
    <div className="pagination">
      <button disabled={page <= 1} onClick={() => onPageChange(page - 1)}>
        Prev
      </button>

      {pageNumbers.map((num) => (
        <button
          key={num}
          className={num === page ? "active" : ""}
          onClick={() => onPageChange(num)}
        >
          {num}
        </button>
      ))}

      <button
        disabled={page >= totalPages}
        onClick={() => onPageChange(page + 1)}
      >
        Next
      </button>
    </div>
  );
}

export default Pagination;
\end{lstlisting}

\begin{lstlisting}[language=JavaScript]
// frontend/src/pages/TaskList.jsx
import { useState, useEffect } from "react";
import api from "../services/api";
import Pagination from "../components/Pagination";

function TaskList() {
  const [tasks, setTasks] = useState([]);
  const [page, setPage] = useState(1);
  const [totalPages, setTotalPages] = useState(1);

  useEffect(() => {
    const load = async () => {
      const { data } = await api.get(`/tasks?page=${page}&limit=10`);
      setTasks(data.data);
      setTotalPages(data.totalPages);
    };
    load();
  }, [page]);

  return (
    <div>
      {tasks.map((task) => (
        <div key={task._id}>{task.title}</div>
      ))}
      <Pagination page={page} totalPages={totalPages} onPageChange={setPage} />
    </div>
  );
}

export default TaskList;
\end{lstlisting}

\begin{tipbox}[TIP]
Keep \texttt{page} in the URL's query string (e.g.\ via
\texttt{useSearchParams} from React Router) so a reload or a shared link
lands on the same page instead of always resetting to page 1.
\end{tipbox}
